
Indeed, \(S\) is affine, so by the first assertion of
\SGARefLink{sga3:xxii:cor:5.9.7}{Exp.~XXII, 5.9.7}, if \(G\) has a Borel subgroup \(B\), it also has a
Killing couple \(B\supset T\). Then, because
\(\operatorname{Pic}(D)=0\), the module \(\mathfrak g_D\) has a
nowhere-zero section \(X\).

Continue to let \(A\) be a semilocal ring and
\(S=\operatorname{Spec}(A)\). For every reductive \(S\)-group \(G\),
the morphism
\[
  \operatorname{Bor}(G)\longrightarrow S
\]
is surjective (because \(G_s\) has Borel subgroups for every
\(s\in S\)), smooth, and projective (\SGARefLink{sga3:xxii:cor:5.8.3}{Exp.~XXII, 5.8.3}). It therefore
has sections after a finite, surjective, étale extension of the
base.\footnote{Editor's note: The original referred to EGA~IV,
\S24, which did not appear. The point is to use ``Bertini's
theorem.'' We give the details of the argument, which were indicated
to us by O.\ Gabber. Let \(X\to S\) be a surjective smooth projective
morphism. Replacing \(S\) by a connected component, one may assume
that \(X/S\) has relative dimension \(d\) at every point
(cf.\ EGA~IV\textsubscript{4}, 17.10.2). One may also assume that
\(X\) is a closed subscheme of a projective space
\(\mathbb P^n_S=\mathbb P(E)\), where \(E\) is a free \(A\)-module
of rank \(n+1\) (cf.\ EGA~II, 5.3.3). Let
\(s_1,\ldots,s_r\) be the closed points of \(S\). By Bertini's
theorem (see, for example, \textup{[Jou83]}, I, 6.10), there is an
open subset \(U\) of
\[
  P=\mathbb P(E^\ast)^d
\]
with nonempty fibers such that, for every point
\(u=(f_1,\ldots,f_d)\) of \(U_{s_i}\), the intersection of
\(X_{\kappa(u)}\) with the \(d\) hyperplanes of
\(\mathbb P^n_{\kappa(u)}\) defined by \(u\) is étale over
\(\kappa(u)\). After shrinking \(U\), one may assume that \(U\) is
the complement in affine space \(\mathbb A^{nd}_S\) of the zero
locus \(V(Q)\) of a polynomial \(Q\) of degree \(m>0\). An easy
induction on the number of variables then shows that \(V(Q)\) cannot
contain every rational point of
\(\mathbb A^{nd}_{\kappa(s_i)}\) if
\(\lvert\kappa(s_i)\rvert>m\). Since the morphism
\[
  A\longrightarrow\prod_i\kappa(s_i)
\]
is surjective, one can find a monic polynomial \(R\in A[X]\) of
degree \(m\) whose image in \(\kappa(s_i)[X]\) is irreducible when
\(\kappa(s_i)\) is finite and has \(m\) distinct roots when
\(\kappa(s_i)\) is infinite. Put
\[
  A'=A[X]/(R),
  \qquad
  S'=\operatorname{Spec}(A').
\]
Then \(S'\to S\) is étale, finite, and surjective, and the residue
field at each of its closed points \(s'_1,\ldots,s'_t\) has
cardinality at least \(2m\).
Thus \(U\) has a rational point \(u_i\) above every closed point of
\(S'\), and, since
\[
  A'\longrightarrow\prod_i\kappa(s'_i)
\]
is surjective \emph{(Translator's note: the French text prints
\(\kappa(s_i)\); the intended closed points of \(S'\) are
\(s'_i\))}, these points lift to a section \(u\) of \(P_{S'}\).
Let \(Z\) be the intersection of \(X_{S'}\) with the \(d\)
hyperplanes of \(\mathbb P^n_{S'}\) defined by \(u\), and let \(V\)
be the open subset of \(Z\) consisting of the points at which \(Z\)
is étale over \(S'\). By EGA~IV\textsubscript{3}, 11.3.8, \(V\)
contains the fibers \(Z_{s'_i}\) for every \(i\). Since
\(\pi:Z\to S'\) is proper, it follows that the closed set
\(\pi(Z-V)\) is empty, whence \(V=Z\). Thus \(Z\to S'\) is
surjective, étale, and proper, hence finite, as is the composite
\(Z\to S\); this gives the desired section of \(X\to S\).}
It follows that:

\SGARefTarget{sga3:xxiv:3.9.2}
\medskip\noindent\textbf{Corollary 3.9.2.} ---
{\itshape Let \(S\) be a semilocal scheme and \(G\) a reductive
\(S\)-group. There is an étale, finite, and surjective morphism
\(S'\to S\) such that \(G_{S'}\) is quasi-splittable.}

\SGARefTarget{sga3:xxiv:3.9.3}
\medskip\noindent\textbf{Remark 3.9.3.} ---
Under the preceding hypotheses, let \(T\) be a maximal torus of
\(G\) (cf.\ \SGARefLink{sga3:xiv:3.20}{Exp.~XIV, 3.20}). One may moreover assume that \(G_{S'}\)
is quasi-splittable relative to \(T_{S'}\): it suffices to apply the
preceding argument to the ``scheme of Borel subgroups containing
\(T\),'' which is finite and étale over \(S\)
(\SGARefLink{sga3:xxii:corollary:5-5-5:item:ii}{Exp.~XXII, 5.5.5(ii)}).

\SGARefTarget{sga3:xxiv:3.10}
\medskip\noindent\textbf{3.10.} ---
Let \(E\) and \(E'\) be two quasi-pinnings of the reductive
\(S\)-group \(G\). There is a unique inner automorphism of \(G\)
carrying \(E\) to \(E'\). Indeed, one reduces immediately to the
split case, where the assertion has already been proved (\SGARefLink{sga3:xxiv:1.5}{1.5}; it
suffices to observe that for an inner automorphism of \(G\), preserving
a pinning is equivalent to preserving the underlying quasi-pinning).
One concludes, as in \S1, that a


quasi-pinning of the reductive \(S\)-group \(G\) defines a splitting
\(h\) of the exact sequence:

\[
  \begin{tikzcd}[column sep=small]
1 \arrow[r]
&
\underline{\operatorname{ad}}(G) \arrow[r]
&
\underline{\Aut}_{S\text{-}\mathrm{gr.}}(G)
  \arrow[r,shift right=.7ex,"p"']
&
\underline{\Autext}(G)
  \arrow[l,shift right=.7ex,dotted,"h"']
  \arrow[r]
&
1
\end{tikzcd}

\]

the image of \(h\) being the subgroup of
\(\underline{\Aut}_{S\text{-}\mathrm{gr}}(G)\) that preserves the
quasi-pinning.

Likewise, if \(G\) and \(G'\) are two quasi-pinned \(S\)-groups, one
defines naturally the subfunctor
\[
  \underline{\operatorname{Isom}}_{S\text{-}\mathrm{gr},\,q\text{-}\acute{e}p}
    (G,G')
\]
of
\(\underline{\operatorname{Isom}}_{S\text{-}\mathrm{gr}}(G,G')\).
The projection from
\(\underline{\operatorname{Isom}}_{S\text{-}\mathrm{gr}}(G,G')\)
to \(\underline{\operatorname{Isomext}}(G,G')\) induces an
isomorphism
\[
  (*)\qquad
  \underline{\operatorname{Isom}}_{S\text{-}\mathrm{gr},\,q\text{-}\acute{e}p}
    (G,G')
  \xrightarrow{\ \sim\ }
  \underline{\operatorname{Isomext}}(G,G').
\]

\SGARefTarget{sga3:xxiv:3.11}
\medskip\noindent\textbf{Theorem 3.11.} ---
{\itshape Let \(S\) be a scheme, let \(\mathscr R\) be a pinned reduced
root datum such that
\(\operatorname{\acute{E}p}_S(\mathscr R)\) exists
(cf.\ Exp.~XXV), and let \(E\) be its group of automorphisms.
Consider the following three categories:
\begin{enumerate}
\item[\SGAInlineRefTarget{sga3:xxiv:3.11:item:i}(i)] the category \(\mathsf{Rev}\) of principal Galois coverings
of \(S\) with group \(E\) (the morphisms are isomorphisms of
\(E\)-bundles);
\item[\SGAInlineRefTarget{sga3:xxiv:3.11:item:ii}(ii)] the category \(\mathsf{Redext}\), whose objects are the
reductive \(S\)-groups of type \(\mathscr R\) (\SGARefLink{sga3:xxii:def:2.7}{Exp.~XXII, 2.7}), the
morphisms from \(G\) to \(G'\) being the elements of
\(\underline{\operatorname{Isomext}}(G,G')(S)\);
\item[(iii)] the category \(\mathsf{Q\acute{e}p}\) of quasi-pinned
reductive \(S\)-groups of type \(\mathscr R\) (the morphisms are the
isomorphisms preserving the quasi-pinnings).
\end{enumerate}
These three categories are equivalent. More precisely, there is a
diagram of functors, commutative up to functorial
isomorphisms:}

\[
  \begin{tikzcd}[column sep=large,row sep=large]
\mathsf{Rev}
  \arrow[rr,"q\acute{e}p"]
&&
\mathsf{Q\acute{e}p}
  \arrow[dl,"i"]
\\
&
\mathsf{Redext}
  \arrow[ul,"\operatorname{rev}"]
&
\end{tikzcd}

\]

We describe these three functors below, leaving to the reader the
verification that the diagram commutes.\footnote{Editor's note: In
\SGARefLink{sga3:xxiv:3.11.4}{3.11.4} we have sketched the verification that the functors \(i\),
\(\operatorname{rev}\), and \(q\acute{e}p\) are fully faithful and
that the composite
\(\operatorname{rev}\circ i\circ q\acute{e}p\) is isomorphic to the
identity functor of \(\mathsf{Rev}\).}

\SGARefTarget{sga3:xxiv:3.11.1}
\medskip\noindent\textbf{3.11.1. The functor \(i\).} ---
This is the evident functor:
\[
  i(G)=G,
  \qquad
  i(f)=\text{the image of \(f\) under the isomorphism 3.10 \((*)\)}.
\]

\SGARefTarget{sga3:xxiv:3.11.2}
\medskip\noindent\textbf{3.11.2. The functor
\(q\acute{e}p\).} ---
Let \(S'\) be a principal Galois covering of \(S\) with group \(E\).
Let
\[
  \bigl(G,T,M,\mathscr R,\Delta,
    (X_\alpha)_{\alpha\in\Delta}\bigr)
\]
be a pinning of
\(G=\operatorname{\acute{E}p}_S(\mathscr R)\), and let
\((X'_\alpha)_{\alpha\in\Delta}\) be the corresponding pinning of
\[
  G'=G_{S'}=\operatorname{\acute{E}p}_{S'}(\mathscr R).
\]
By \SGARefLink{sga3:XXIII:5.5bis}{Exp.~XXIII, 5.5 bis}, there is a unique group homomorphism
\[
  \theta:E
  \longrightarrow
  \underline{\Aut}_{S\text{-}\mathrm{gr}}
    \bigl(\operatorname{\acute{E}p}_S(\mathscr R)\bigr)
  =A(\mathscr R)(S)
\]


such that, for every \(h\in E\), the automorphism \(\theta(h)\)
induces \(D_S(h)\) on \(T\), and
\[
  \Lie(\theta(h))(X_\alpha)=X_{h(\alpha)}
  \qquad(\alpha\in\Delta).
\]
Taking into account the action of \(E\) on \(S'\), this yields an
action of \(E\) on \(G'\), compatible with the action of \(E\) on
\(S'\) and preserving the canonical quasi-pinning of
\(G'\).\footnote{Editor's note: The preceding discussion expands the
original in order to show that the action of \(E\) on \(G'\) is
obtained by combining the natural action of \(E\) on
\(\operatorname{\acute{E}p}_{S'}(\mathscr R)\) with the given action
on \(S'\). The canonical Killing couple of \(G'\) is preserved by
this action, as is the quasi-pinning given by the element
\(\widetilde X'\) of
\[
  \Gamma\!\left(
    \Delta_{S'},
    \Lie(G'/S')_{\Delta_{S'}}
  \right),
\]
equal to \(X'_\alpha\) on the copy of \(S'\) indexed by \(\alpha\).
Indeed, since \(E\) acts on \(\Delta_{S'}\) by permuting the copies
of \(S'\), one has \(h(\widetilde X')=\widetilde X'\) for every
\(h\in E\).}
Since \(S'\to S\) is a covering for the \textup{(fpqc)} topology, it
is an effective descent morphism for the fibered category of affine
morphisms. Write
\[
  q\acute{e}p(S')
  =
  Q\text{-}\operatorname{\acute{E}p}_{S'/S}(\mathscr R)
\]
for the quasi-pinned \(S\)-group obtained by Galois
descent.\footnote{Editor's note: See, for example, TDTE~I, p.~22,
\SGARefLink{sga3:xxiv:section:1}{Example~1}. One may also describe \(q\acute{e}p(S')\) as the twist of
\(G=\operatorname{\acute{E}p}_S(\mathscr R)\) by the \(E\)-torsor
\(S'/S\), that is, the \textup{(fpqc)} sheaf quotient of
\(S'\times_SG\) by the right action of \(E\) defined by
\[
  (s',g)\cdot h=(s'h,h^{-1}(g)).
\]
Since \(G\) is affine over \(S\), and since \(E\) acts on \(G\) by
group automorphisms, this sheaf is represented by an \(S\)-group
\(G^\sharp\), which is a ``twisted'' form of \(G\), and
\(D^\sharp=\operatorname{Dyn}(G^\sharp)\) is the twist of the
constant Dynkin scheme \(\Delta_S\) by the torsor \(S'/S\). Since
\(E\) normalizes \(B\) and \(T\), one obtains likewise a Killing
couple \((B^\sharp,T^\sharp)\) of \(G^\sharp\). On the other hand,
put
\(\mathfrak g=\Lie(G/S)\) and
\(\mathfrak g^\sharp=\Lie(G^\sharp/S)\). For every \(U\to S\), the
sections of \(\mathfrak g^\sharp\) over \(U\) are the
\(E\)-equivariant \(S\)-morphisms
\[
  U\times_SS'\longrightarrow W(\mathfrak g).
\]
Since \(D^\sharp\times_SS'\) is \(E\)-isomorphic to
\(\Delta\times S'\), endowed with the action
\[
  (\alpha,s')\cdot h=(h^{-1}(\alpha),s'h),
\]
the morphism \((\alpha,s')\mapsto X_\alpha\) gives a section of
\(\mathfrak g^\sharp\) over \(D^\sharp\) which is a quasi-pinning,
that is, a nowhere-zero section of
\[
  \bigl(
    \mathfrak g^\sharp\otimes_{\mathcal O_S}\mathcal O_{D^\sharp}
  \bigr)^{D^\sharp}
\]
(cf.\ \SGARefLink{sga3:xxiv:3.8}{3.8}).
}

\SGARefTarget{sga3:xxiv:3.11.3}
\medskip\noindent\textbf{3.11.3. The functor
\(\operatorname{rev}\).} ---
Let \(G\) be a reductive \(S\)-group of type \(\mathscr R\). Write
\[
  \operatorname{rev}(G)
  =
  \underline{\operatorname{Isomext}}
    \bigl(\operatorname{\acute{E}p}_S(\mathscr R),G\bigr).
\]
This is a principal homogeneous bundle under \(E_S\)
(cf.\ \SGARefLink{sga3:xxiv:1.10}{1.10} and \SGARefLink{sga3:xxiv:1.19}{1.19}), that is, an object of \(\mathsf{Rev}\).

\SGARefTarget{sga3:xxiv:3.11.4}
\medskip\noindent\textbf{3.11.4.} ---
It is clear from the isomorphism \SGARefLink{sga3:xxiv:3.10}{3.10} \((*)\) that the functor \(i\)
is fully faithful. On the other hand, if \(G_0,G,G'\) are reductive
\(S\)-groups, the natural morphism
\[
  \underline{\operatorname{Isom}}(G_0,G)
  \times_S
  \underline{\operatorname{Isom}}(G,G')
  \longrightarrow
  \underline{\operatorname{Isom}}(G_0,G')
\]
induces a morphism
\[
  (\star)\qquad
  \underline{\operatorname{Isomext}}(G_0,G)
  \times_S
  \underline{\operatorname{Isomext}}(G,G')
  \longrightarrow
  \underline{\operatorname{Isomext}}(G_0,G'),
\]
since, for every \(S'\to S\), if
\[
  g_0\in G_0(S'),\quad g\in G(S'),\quad
  u\in\underline{\operatorname{Isom}}(G_0,G)(S'),\quad
  v\in\underline{\operatorname{Isom}}(G,G')(S'),
\]
one has
\[
  v\circ\operatorname{int}(g)\circ u\circ\operatorname{int}(g_0)
  =
  v\circ u\circ
  \operatorname{int}\bigl(u^{-1}(g)g_0\bigr).
\]
Taking \(G_0=\operatorname{\acute{E}p}_S(\mathscr R)\), one obtains
a map
\[
\begin{aligned}
  \underline{\operatorname{Isomext}}(G,G')(S)
  \longrightarrow
  \operatorname{Hom}_{\mathsf{Rev}}\bigl(
  &\underline{\operatorname{Isomext}}
    (\operatorname{\acute{E}p}_S(\mathscr R),G),\\
  &\underline{\operatorname{Isomext}}
    (\operatorname{\acute{E}p}_S(\mathscr R),G')
  \bigr),
\end{aligned}
\]
and, to show that this map is a bijection, one may assume by
\textup{(fpqc)} descent that \(G\) and \(G'\) are pinned, in which
case it is evident.


Similarly, if \(S_1,S_2\) are two objects of \(\mathsf{Rev}\), and
if one writes
\[
  S_i\times^E\operatorname{\acute{E}p}_S(\mathscr R)
\]
for the \(S\)-group \(q\acute{e}p(S_i)\) obtained by twisting
\(\operatorname{\acute{E}p}_S(\mathscr R)\) by the \(E\)-torsor
\(S_i\), there is a natural map
\[
\begin{aligned}
  \operatorname{Hom}_{\mathsf{Rev}}(S_1,S_2)
  \longrightarrow
  \operatorname{Hom}_{\mathsf{Q\acute{e}p}}\bigl(
    &S_1\times^E\operatorname{\acute{E}p}_S(\mathscr R),\\
    &S_2\times^E\operatorname{\acute{E}p}_S(\mathscr R)
  \bigr).
\end{aligned}
\]
To show that this map is a bijection, one may assume by
\textup{(fpqc)} descent that \(S_1\) and \(S_2\) are trivial
\(E\)-bundles, in which case the assertion is evident.

Thus the three functors in the diagram are fully faithful. Moreover,
for every object \(S'\) of \(\mathsf{Rev}\), there is a natural
morphism
\[
  S'
  \longrightarrow
  \underline{\operatorname{Isomext}}\!\left(
    \operatorname{\acute{E}p}_S(\mathscr R),
    S'\times^E\operatorname{\acute{E}p}_S(\mathscr R)
  \right),
\]
and to see that it is an isomorphism, one may assume by
\textup{(fpqc)} descent that \(S'\) is a trivial \(E\)-bundle, in
which case the assertion is evident.

There is therefore an isomorphism of functors
\[
  \operatorname{Id}_{\mathsf{Rev}}
  \longrightarrow
  \operatorname{rev}\circ i\circ q\acute{e}p.
\]
Since \(\operatorname{rev}\) is fully faithful, for every reductive
\(S\)-group \(H\) (not necessarily quasi-split) one obtains a
bijection
\[
  \underline{\operatorname{Isomext}}\!\left(
    Q\text{-}\operatorname{\acute{E}p}_{\operatorname{rev}(H)/S}
      (\mathscr R),
    H
  \right)(S)
  \simeq
  \operatorname{Aut}_{\mathsf{Rev}}\bigl(\operatorname{rev}(H)\bigr),
\]
functorial in \(H\). In particular, the identity morphism of
\(\operatorname{rev}(H)\) corresponds to a ``canonical'' element
\[
  u_0\in
  \underline{\operatorname{Isomext}}\!\left(
    Q\text{-}\operatorname{\acute{E}p}_{\operatorname{rev}(H)/S}
      (\mathscr R),
    H
  \right)(S).
\]
Moreover, every element \(u\) of this same set corresponds to an
automorphism \(\varphi_u\) of \(\operatorname{rev}(H)\), and
\(q\acute{e}p(\varphi_u)\) is the unique automorphism of the pinned
\(S\)-group
\[
  Q\text{-}\operatorname{\acute{E}p}_{\operatorname{rev}(H)/S}
    (\mathscr R)
\]
such that
\[
  u_0\circ q\acute{e}p(\varphi_u)=u.
\]
This proves \SGARefLink{sga3:xxiv:3.11}{Theorem 3.11}.

We develop one of the corollaries of \SGARefLink{sga3:xxiv:3.11}{3.11}.

\SGARefTarget{sga3:xxiv:3.12}
\medskip\noindent\textbf{Corollary 3.12.} ---
{\itshape For every reductive \(S\)-group \(G\), there is a quasi-pinned
\(S\)-group \(G_{q\text{-}\acute{e}p}\) and an ``outer isomorphism''}
\[
  u\in
  \underline{\operatorname{Isomext}}
    (G_{q\text{-}\acute{e}p},G)(S).
\]
{\itshape The pair \((G_{q\text{-}\acute{e}p},u)\) is unique up to a
unique isomorphism.}

Indeed, one may assume that \(G\) has constant type \(\mathscr R\),
and one takes
\[
  G_{q\text{-}\acute{e}p}
  =
  Q\text{-}\operatorname{\acute{E}p}_{\operatorname{rev}(G)/S}
    (\mathscr R).
\]

\SGARefTarget{sga3:xxiv:3.12.1}
\medskip\noindent\textbf{3.12.1.} ---
Observe that the data of the pair
\((G_{q\text{-}\acute{e}p},u)\) canonically define the \(S\)-scheme
(cf.\ \SGARefLink{sga3:xxiv:1.11}{1.11})
\[
  Q=
  \underline{\operatorname{Isomint}}
    (G_{q\text{-}\acute{e}p},G),
\]
which is a principal homogeneous bundle under
\(\operatorname{ad}(G_{q\text{-}\acute{e}p})\), and whose data are
``equivalent'' to those of the inner twisted form \(G\) of
\(G_{q\text{-}\acute{e}p}\). Moreover, \(Q\) is precisely the
``scheme of quasi-pinnings of \(G\),'' a definition that explains its
structure as a principal homogeneous bundle (on the left) under
\(\operatorname{ad}(G)\) (\SGARefLink{sga3:xxiv:3.10}{3.10}); compare \SGARefLink{sga3:xxiv:1.20}{1.20}.
